\newpage

\begin{center}
    {\Large\itshape Alban à l'heure du troubadour « Pierre d'Alban »}
\end{center}

\vspace{1.2em}

\lettrine[lines=3,lhang=0.1,nindent=0em]{N}{otre} poète ambulant n'eut que l'embarras du choix pour faire retentir sa chanson dans les « castels » du voisinage. Sa devise « audenti succedit opus » -- l'audace paie -- dut l'amener au château de Paulin, à celui de Miolles où les chevaliers de Matha devaient s'illustrer par un des leurs, Saint Jean de Matha. La Tapie dut avoir sa visite. Un sieur de Pujol au XIV\textsuperscript{e} siècle y sera qualifié de « plus fort pétardier de France ». Il est probable qu'il fut accueilli à Massals où les Brandouin ont laissé gravés dans la pierre les panaches de leurs armes. Les châteaux de Nougaret, d'Alban, de Saint-André, de Coupiac durent l'accueillir. S'il s'aventura sur les berges du Tarn, il fut peut-être l'hôte des « de Sales » dont un cousin devient évêque de Genève, le célèbre Saint François de Sales. Poussa-t-il jusqu'à la demeure du vicomte d'Ambialet et chez les moines bénédictins du prieuré voisin ? C'est vraisemblable. Quel que soit le seigneur de haut lignage, quelle que fut la dame « de haute cuisse » qui ait allumé cette étoile filante au ciel de l'Albigeois, nous avons la chance qu'un écho de la voix de charme de notre troubadour ait récemment retenti grâce à Annick Roche dont voici le récit\ldots
